\sscp{\tsc{Babar}}{04-06-05}
There is strong evidence for a \tsc{Babar} node.
This group includes all the languages of the \ili{Babar} archipelago
(except \ili{Wetan} which is a member of the \ili{Luang} language/dialect
cluster), as well as West \ili{Damer} [drn] on \ili{Damer} island to the north of \ili{Babar}.
Within \tsc{Babar}, we can identify a \tsc{Peripheral Babar}
node (\srf{04-06-05-01}) and a \tsc{South Babar} node (\srf{04-06-05-02}).
The languages of the \tsc{Babar} node are poorly documented.
The 193-item word lists in \citet{Taber1993} are the only source
of data for five of these languages:\footnote{
		During the final stages of preparation of
		this book for publication, Edwards
		collected additional data on all \tsc{South Babar}
		languages, as well as the North \ili{Babar} language.
		This data has been archived as \citet{Edwards2023}.
		With the exception of \ili{Emplawas} and West \ili{Marsela},
		this data has not been integrated into this section.}
\ili{Dai} [dij], North \ili{Babar} [bcd], East \ili{Marsela} [vme], \ili{Imroing} [imr], and \ili{Tela-Masbuar} [tvm].\footnote{
		We usually
		refer to \ili{Dawera-Daweloor} as \it{Dawera}
		and \ili{Tela-Masbuar} as \it{Tela} as these are the locations for
		which data is available.
		Both \it{Masela} and \it{Marsela} are used for the island to
		the south of \ili{Babar}.
		We follow \citeauthor{vanDijk2019} in using \it{Marsela},
		which accords with local usage.}

Additional data comes from
the following sources: \citet{Steinhauer2009} for South\ili{east} \ili{Babar}
[vbb], \citet{Chlenova2002,Chlenova2012} for \ili{Dawera-Daweloor} [ddw],
\citet{ChlenovChlenova2008b} for West \ili{Damer},
\citet{Edwards2023} for \ili{Emplawas} and West \ili{Marsela}, as well as \citet{vanDijk2021}
			for \ili{Central} \ili{Marsela}, and \ili{Serili}.\footnote{
			The data in \citet{vanDijk2021}
			was collected between 1981--1983.}
\citet{Zobel2018} discusses
the phonological history of \tsc{South Babar},
based on the data in \citet{Taber1993},
and also provides a list of Proto-South \ili{Babar} reconstructions.
Our interpretation of the data differs in
several respects to that of \citet{Zobel2018}.

Throughout this section data from \citet{Taber1993}
has been re-transcribed according to a preliminary
phonological interpretation.
Data from other sources follows that source,
with the exception of mid-vowels
for which we use the symbols in \citet{Taber1993} wherever possible.\footnote{
		Where both \citet{Taber1993} and another source have
		a cognate word for the same language, we use the vowel symbols
		from \citet{Taber1993}. For instance, \citet[418]{Taber1993}
		gives \ili{Dawera-Daweloor} ‹ˈmɛtmel› `black' while \citet[165]{Chlenova2002}
		gives ‹metmel› `black'. Thus, we transcribe this word \it{mɛtmel}.}
This is because there is evidence that at l\ili{east} some languages
of the \ili{Babar} Islands have a contrast between mid-high and mid-low
vowels\footnote{
		Edwards' fieldwork in mid 2025 indicates
		that all members of the \tsc{South Babar} node
		have a contrast between mid-high and mid-low front vowels (/e/ vs. /ɛ/).
		Some also have a contrast between two mid back vowels (/o/ vs. /ɔ/).
		Minimal pairs from \ili{Emplawas} include: /we/ [β̞e] ‘water’ vs.
		/wɛ/ [β̞ɛ] ‘rope’ %, /k-ep/ ‘\tsc{1pl.incl}-drink’ vs. /kɛp/ ‘Tepa’,
		and /kɛ-wel/ [kɛβ̞el] ‘\tsc{1pl.incl}-cut’
		vs. /kɛ-wɛl/ [kɛβ̞ɛl] ‘\tsc{1pl.incl}-buy’.}
which may be under-represented in the available sources.\footnote{
		The data in \citet{Chlenova2002} for Dawera and
		\citet{ChlenovChlenova2008b} for West \ili{Damer}
		was provided by their consultants in written form.
		The West \ili{Marsela} and \ili{Serili} data in \citet{vanDijk2021}
		was similarly compiled by native-speaker consultants,
		while the \ili{Central} \ili{Marsela} data was compiled by \citeauthor{vanDijk2021}
		with help from her consultants.
		None of these sources indicates a distinction
		between mid-high and mid-low vowels.
		Edwards' fieldwork confirmed a contrast
		between mid /e o/ and mid-low /ɛ ɔ/ in \ili{Serili},
		as well as mid-high /e/ and mid-low /ɛ/
		in West \ili{Marsela} and \ili{Central} \ili{Marsela}.
		While \citet{Steinhauer2009} does not report a contrast
		between mid-vowel heights for SE. \ili{Babar},
		Edwards' fieldwork found such a contrast
		for some SE. \ili{Babar}-speaking villages.}

The \tsc{Babar} node is primarily defined by fortition of PMP *s.
The \tsc{Peripheral Babar} languages have *s > \it{d},
and the \tsc{South Babar} languages have *s > \it{t}.
We propose that these reflexes are the result of an intermediate retroflex **ʈ.
Evidence in favour of medial **ʈ is discussed below.
The \tsc{Peripheral Babar} languages have sporadic *s > \it{t}
in syllable codas and before voiceless consonants,
though such instances of [t] could be allophones of /d/.
Examples of *s > **ʈ are given in \trf{04-45}.
(See also reflexes of *ma-nipis ‘thin’ in \trf{04-01} on \prf{tab:04-01}.)

The change of PMP *s > \tsc{Peripheral Babar} \it{d} is extremely
unusual and requires discussion.
This change must be proposed instead of more common -- though
still rare---*s > **t as the reflexes of PMP *s
and *t remain distinct.
\tsc{South Babar} has PMP *t > **k,
but PMP *s > \it{t},
while other \tsc{Babar} languages have *s > \it{d}, but *t = \it{t}.
The simplest solution is to posit *s > **ʈ at Proto-\ili{Babar},
with *t > **k occurring before **ʈ > \it{t} in Proto-South \ili{Babar}.


\begin{table}\tcp{\tsc{Babar} *s > **ʈ}{04-45}
\begin{adjustbox}{width=\textwidth}
	\begin{threeparttable}\st{2pt}
	\begin{tabular}{ll@{}ll@{}ll@{}l}\lsptoprule
local gl.	&	{\hr}br\ili{east}											&	{\hr}who							&	{\hr}one	&	{\hr}dog	&	{\hr}smoke	&	{\hr}teeth\\
PMP				&	\hr*\tbr{s}u\tbr{s}u						&	\hr*\tbr{s}ai					&	\hr*ma-ə\tbr{s}a	&	\hr*a\tbr{s}u	&	**ma-qa\tbr{s}u	&	\hr*ŋi\tbr{s}i\su{Δ}\\
P.-\ili{Babar}	&	**\tbr{ʈ}u\tbr{ʈ}u							&	**\tbr{ʈ}ai						&	**me\tbr{ʈ}(en)	&	**a\tbr{ʈ}u	&	**ma\tbr{ʈ}u	&	**ŋi\tbr{ʈ}i\\ \midrule
W. \ili{Damer}	&	{\hr\hr}\tbr{d}u\tbr{t}h-o			&	{\hr\hr}\tbr{d}esɛri	&	\hr\cg{(mɛhn-o)}\su{†}	&	{\hr\hr}o\tbr{t}h-o	&	\tcg{(o muhal-o)}\su{‡}	&	\hr\cg{(nin-o)}\\
Dawera		&	{\hr\hr}\tbr{d}u\tbr{d}k-ol			&	{\hr\hr}\tbr{d}eyɛ		&	{\hr\hr}mɛ\tbr{d}	&	{\hr\hr}a\tbr{d}k-ol	&	am mɛ\tbr{d}iky	&	{\hr\hr}i\tbr{d}-ol\\
\ili{Dai}				&	{\hr\hr}\tbr{d}u\tbr{d}ʔ-on			&	{\hr\hr}\tbr{d}eyehe	&	{\hr\hr}mɛ\tbr{d}	&	{\hr\hr}a\tbr{d}y-on	&	ma\tbr{d}ʔ-ɛn-on	&	{\hr\hr}i\tbr{d}-on\\
N. \ili{Babar}	&	{\hr\hr}\tbr{d}u\tbr{t}-nie			&	{\hr\hr}\tbr{d}ediɛ	&	{\hr\hr}mɛ\tbr{d}e	&	{\hr\hr}a\tbr{d}ui-ai	&	am-e ma\tbr{t}-ni-ai	&	{\hr\hr}i\tbr{t}-ne\\
C. \ili{Marsela}&	{\hr\hr}\tbr{t}u\tbr{t}-ni			&	{\hr\hr}\tbr{t}yoi	&	{\hr\hr}mɛ\tbr{t}	&	{\hr\hr}e\tbr{t}y	&	\tcg{(owour i-ei)}	&	\hr\cg{(ninn)}\\
W. \ili{Marsela}&	{\hr\hr}\tbr{t}u\tbr{t}y				&	{\hr\hr}\tbr{t}io i	&	{\hr\hr}mɛ\tbr{t}n	&	{\hr\hr}u\tbr{t}y	&	\tcg{(owr)}	&	\hr\cg{(nin)}\\
E. \ili{Marsela}&	{\hr\hr}\tbr{t}u\tbr{t}					&	{\hr\hr}\tbr{t}oi	&	{\hr\hr}mɛ\tbr{t}el	&	{\hr\hr}u\tbr{t}	&	ui mo\tbr{t}y	&	\hr\cg{(lil)}\\
\ili{Serili}		&	{\hr\hr}\tbr{t}u\tbr{t}					&	{\hr\hr}\tbr{t}oi	&	{\hr\hr}mɛ\tbr{t}el	&	{\hr\hr}u\tbr{t}	&	\tcg{(ui-ho bo{\tl}bor-ɛ)}	&	\hr\cg{(lil)}\\
SE. \ili{Babar}	&	{\hr\hr}\tbr{t}u\tbr{t}					&	{\hr\hr}\tbr{t}oi	&	{\hr\hr}me\tbr{t}el	&	{\hr\hr}u\tbr{t}y	&	ui mo{\tl}mo\tbr{t}	&	\hr\cg{(lil)}\\
\ili{Emplawas}	&	{\hr\hr}\tbr{t}u\tbr{t}y-o			&	{\hr\hr}\tbr{t}oi	&	{\hr\hr}mɛ\tbr{t}	&	{\hr\hr}u\tbr{t}y-o	&	mu\tbr{t}y-o	&	\hr\cg{(lil-wo)}\\
\ili{Imroing}		&	{\hr\hr}\tbr{t}u\tbr{t}u-mu-te	&	{\hr\hr}\tbr{t}oi	&	{\hr\hr}mɛ\tbr{t}le	&	{\hr\hr}u\tbr{t}i-e	&	mu\tbr{t}i-le	&	\hr\cg{(lil-mu-te)}\\
Tela			&	{\hr\hr}\tbr{t}u\tbr{t}-ni			&	{\hr\hr}\tbr{t}eʔoi	&	{\hr\hr}me\tbr{t}nu	&	{\hr\hr}nu\tbr{t}i-e	&	\tcg{(obri-e)}	&	\hr\cg{(nin-e)}\\
		\lspbottomrule
	\end{tabular}
		\begin{tablenotes}\tnsize
			\item[†]	{West \ili{Damer} \it{mɛhn-o} ‘one’ is probably borrowed from East \ili{Damer} \it{mehan} ‘one’.}
			\item[‡]	{West \ili{Damer} \it{muhal-o} ‘smoke’ is probably a borrowing.
								Compare East \ili{Damer} \it{muklan},
								\ili{Luang} \it{moklə}, \ili{Wetan} \it{mokla} ‘smoke’,
								any of which are possible sources.}
			\item[Δ]	{All words in brackets in the *ŋisi column are reflexes of *nipən.}
		\end{tablenotes}
	\end{threeparttable}
\end{adjustbox}
\end{table}

How the change of *s > \it{d} in \tsc{Peripheral Babar} came
about is illuminated by a regional perspective on the phonetics of coronal plosives.
It is common in the Greater Timor region for the voiceless apical/laminal
plosive /t/ to be phonetically dental [t̪] while the voiced apical/laminal
plosive /d/ is phonetically alveolar [d̳] or even slightly retroflex [ḍ].
Similarly, the voiceless fricative /s/ is alveolar [s̳].
Some languages of the southwest Maluku area have taken this difference
in place of articulation a step further with what was (historically)
a voiced alveolar plosive /d/ realised as fully retroflex.
This is clearly exemplified by \ili{Kisar} which contrasts fully retroflex
(and voiceless) /ʈ/ with non-retroflex /t/ \citep{ChristensenChristensen1992}.
\ili{Kisar} /ʈ/ is from PMP *d
and is orthographically represented with the letter ‹d›.

This suggests that the change of *s > \it{d} in \tsc{Peripheral
Babar} went through an intermediate voiceless retroflex or apical
stage with this sound subsequently reanalysed as /d/.
Further support for this comes from \ili{Dai} in which voiced [d]
and voiceless retroflex [ʈ] are mostly in complementary environments
in the available data.
They are thus probably analysable as allophones of a single phoneme.
\ili{Dai} [ʈ] mostly occurs in syllable codas while [d] occurs elsewhere.
Examples include: \it{ady-on} [ˈaʈʰon], ‘dog’ \it{n-nod} [n̩ˈnoʈ]
‘run’, \it{nodnon} [ˈnoʈnon] ‘rattan’, \it{nɛ-wrid} [ˈnɛvriʈ] ‘bathe’,
and \it{n-pud} [n̩ˈpuʈ] ‘rotten’.\footnote{
		These \ili{Dai} words can
		be compared with cognates in Dawera which has voiced [d] in all
		environments: \it{adk-ol} [ˈadkol] ‘dog’,
		\it{lod} [lod] ‘run’, \it{lodel} [ˈlodel] ‘rattan’,
		\it{le-wrid} [ˈleʋrid] ‘bathe’, and \it{pud-ɛl} [ˈpudɛl] ‘rotten’.
		Exceptions in \ili{Dai} where retroflex [ʈ] does not occur in coda
		position or [d] does not occur in onsets include \it{mɛd} [mɛd]
		‘one’ and \it{n-dɛr} [n̩ˈʈɛr] ‘dry’ (Dawera \it{dɛr-el} [ˈdɛrel] ‘dry’).
		There are also occasional instances in which coda [t] in \ili{Dai}
		corresponds to **d, such as \it{dudʔ-on} [ˈdutʔon] ‘br\ili{east}’
		and \it{madʔɛn-on} [ˈmatʰʔɛˈnon] ‘smoke’.}

Changes in addition to *s > **ʈ which support the validity of
the \tsc{Babar} node are *z > **h,
*u > \it{i} /{\gap\#},
merger of *b/*w > **w,
and *R > **y, \it{r}.
Examples of final *u > **i,\footnote{
		West \ili{Damer}, \ili{Dai}, and Dawera usually have fortition
		of final *u for nouns (\srf{04-06-05-01}).
		Other \ili{Babar} languages a suffix \it{-i} which occurs on verbs and adjectives.
		These facts, combined with the paucity of data,
		means there are currently few unambiguous cases of final *u >
				**i in \tsc{Babar}.}
as well as *z > **h are given in \trf{04-46}.\footnote{
		An additional example of *z > **h is *zauq > West \ili{Damer} \it{heyo} ‘far’.
		Other \ili{Babar} languages have non-cognate words for ‘far’.}
Word-final *u > \it{i} is also found in nearby \ili{Wetan},
as well as \ili{Masiwang},
\ili{Yamdena} and the \tsc{Bomberai Tip} languages,
all belonging to \tsc{Seram-Tanimbar-Bomberai}
(\crf{ch:STB}).
PMP *z > **h > {\0} in Dawera
and all \tsc{South Babar} languages except \ili{Serili}.

\begin{table}[p]\centering\tcp{\tsc{Babar} *u > *i /{\gap\#} and *z > **h}{04-46}
%\begin{adjustbox}{width=\textwidth}
	\begin{threeparttable}\st{2pt}
	\begin{tabular}{llll|lll}\lsptoprule
local gl.	&	{\hr}three	&	{\hr}seven	&	{\hr}ten	&	{\hr}path	&	green	&	{\hr}rain\\
PMP	&	\hr*təl\tbr{u}	&	\hr*pit\tbr{u}	&	**but\tbr{u}\su{†}	&	\hr*\tbr{z}alan	&	**mo\tbr{z}o\su{‡}	&	\hr*qu\tbr{z}an\\
P.-\ili{Babar}	&	**tel\tbr{i}	&	**it\tbr{i}	&	**wut\tbr{i}	&	**\tbr{h}alan	&	**mo\tbr{h}o	&	**u\tbr{h}an\\ \midrule
W. \ili{Damer}	&	{\hr\hr}wyɛ-tel\tbr{i}	&	{\hr\hr}wit\tbr{i}	&	{\hr\hr}uswut\tbr{i}	&	{\hr\hr}\tbr{h}oll-o	&	{\hr\hr}mɛh{\tl}mɛ\tbr{h}o	&	{\hr\hr}u\tbr{h}an-o\\
Dawera	&	{\hr\hr}tɛl	&	{\hr\hr}it\tbr{y}	&	{\hr\hr}de-wut	&	{\hr\hr}all-ol	&	\hr\cg{(hiʤau)}	&	{\hr\hr}ul-ol\\
\ili{Dai}	&	{\hr\hr}tɛn	&	{\hr\hr}it\tbr{y}	&	{\hr\hr}d-wut	&	{\hr\hr}\tbr{h}ann-on	&	\hr\cg{(hiʤau)}	&	{\hr\hr}u\tbr{h}n-on\\
N. \ili{Babar}	&	{\hr\hr}kin\tbr{i}	&	{\hr\hr}it\tbr{i}	&	{\hr\hr}de-wut\tbr{i}	&	{\hr\hr}\tbr{h}aʔan-ai	&	\hr\cg{(hiʤau)}	&	{\hr\hr}\tbr{h}ona\\
C. \ili{Marsela}	&	{\hr\hr}wɔk-kɛl	&	{\hr\hr}wɔ-ii	&	{\hr\hr}wuk\tbr{i}	&	{\hr\hr}all	&	{\hr\hr}mo{\tl}moa	&	\hr\cg{(omn)}\\
W. \ili{Marsela}	&	{\hr\hr}wo-kɛl	&	{\hr\hr}wo-ik	&	{\hr\hr}wukk\tbr{i}	&	{\hr\hr}al	&		&	\hr\cg{(omn)}\\
E. \ili{Marsela}	&	{\hr\hr}wɔk-kɛl	&	{\hr\hr}wo-ik	&	{\hr\hr}wuk\tbr{i}	&	{\hr\hr}al	&		&	\hr\cg{(umul)}\\
\ili{Serili}	&	{\hr\hr}bok-kɛl	&	{\hr\hr}boʔot	&	{\hr\hr}buk\tbr{i}	&	{\hr\hr}\tbr{h}al-ɛ	&	{\hr\hr}mo{\tl}m\tbr{h}o	&	\hr\cg{(omen)}\\
SE. \ili{Babar}	&	{\hr\hr}wɔ-kɛl\tbr{y}	&	{\hr\hr}wɔ-ix\tbr{y}	&	{\hr\hr}wuk\tbr{i}	&	{\hr\hr}al	&		&	\hr\cg{(uml)}\\
\ili{Emplawas}	&	{\hr\hr}wo-kɛl\tbr{y}	&	{\hr\hr}wo-ik	&	{\hr\hr}wukk\tbr{i}	&	{\hr\hr}as-o	&	{\hr\hr}mo{\tl}moy-i	&	\hr\cg{(op-o)}\\
\ili{Imroing}	&	{\hr\hr}wo-kel\tbr{i}	&	{\hr\hr}wo-ik\tbr{i}	&	{\hr\hr}wuk\tbr{i}	&	{\hr\hr}ann-e	&	\hr\cg{(mot{\tl}mote)}	&	\hr\cg{(oml-e)}\\
Tela	&	{\hr\hr}wo-ken\tbr{y}	&	{\hr\hr}wo-ik\tbr{y}	&	{\hr\hr}wukk\tbr{i}	&	{\hr\hr}ann-e	&	{\hr\hr}mə{\tl}moi	&	\hr\cg{(omn-e)}\\
		\lspbottomrule
	\end{tabular}
		\begin{tablenotes}\tnsize
			\item[†]	{**butu is reconstructed to PCEMP with the meaning ‘group, bunch, cluster’.
								Its reflexes also mean `ten' in \tsc{\ili{Ambon}-Seram} languages.}
			\item[‡]	{Cognates of **mozo ‘green’ occur widely in other \tsc{Timor-Babar} languages.
								See the note to \trf{04-27} (\prf{tab:04-27}) for additional examples.
								\ili{Imroing} \it{mot{\tl}mote} is borrowed from \ili{Wetan} \it{mota}.}
		\end{tablenotes}
	\end{threeparttable}
%\end{adjustbox}
\end{table}

PMP *R and *j merge in \tsc{Babar}
and then undergo an unconditioned split between **y, {\0} and **r.
When *R/*j were lost,
this was probably through intermediate **y in the pre-\ili{Babar} stage
as evidenced by traces of /y/
and raised vowels in some reflexes.
Examples of *R/*j > **y,
{\0} are given in \trf{04-47}
and examples of *R/*j > **r are given in \trf{04-48}.

\begin{table}\centering\tcp{\tsc{Babar} *R/*j > **y, {\0}}{04-47}
\begin{adjustbox}{width=\textwidth}
	\begin{threeparttable}\st{2pt}
	\begin{tabular}{llllllll}\lsptoprule
local gl.	&	{\hr}house	&	{\hr}blood	&	{\hr}dry in sun	&	{\hr}sun	&	{\hr}ySib	&	{\hr}bathe	&	{\hr}name\\
PMP	&	\hr*\tbr{R}umaq	&	\hr*da\tbr{R}aq	&	\hr*wa\tbr{R}i	&	\hr*qalə\tbr{j}aw	&	\hr*hua\tbr{j}i	&	\hr*di\tbr{R}us	&	\hr*ŋa\tbr{j}an\\
P.-\ili{Babar}	&	**\tbr{y}uma	&	**ra\tbr{y}a	&	**wa\tbr{y}a	&	**la\tbr{y}a	&	**wa\tbr{y}i	&	**(wa)riʈ\su{†}	&	**ŋan\\ \midrule
W. \ili{Damer}	&	{\hr\hr}uma	&	{\hr\hr}ro-o	&		&	{\hr\hr}lawo-ni	&	{\hr\hr}wey-o	&	{\hr\hr}wadir-o	&	{\hr\hr}non-o\\
Dawera	&	{\hr\hr}um-ol	&	{\hr\hr}rai-ol	&	{\hr\hr}me-wael	&	{\hr\hr}ley-ol	&	{\hr\hr}wel-ol	&	{\hr\hr}le-wrid	&	\hr\cg{(wulɛtal-ol)}\\
\ili{Dai}	&	{\hr\hr}um-on	&	{\hr\hr}raʔ-on	&	{\hr\hr}rɛ-viai	&	{\hr\hr}ne-on	&	{\hr\hr}wai-ny-on	&	{\hr\hr}nɛ-wrid	&	\hr\cg{(wnyɛtan)}\\
N. \ili{Babar}	&	{\hr\hr}oma-ai	&	{\hr\hr}raʔa-nai	&	{\hr\hr}nɛ-wiai-ɛna	&	{\hr\hr}neya	&	{\hr\hr}wia-ny-ai	&	{\hr\hr}na-wrod-o	&	{\hr\hr}an-ai\\
C. \ili{Marsela}	&	{\hr\hr}ɛm	&	{\hr\hr}rɛ	&	{\hr\hr}mo-wɛy-ei	&	{\hr\hr}loo	&	{\hr\hr}wei-ni	&	{\hr\hr}ne-itir	&	{\hr\hr}nɛn\\
W. \ili{Marsela}	&	{\hr\hr}im	&	{\hr\hr}rɛ	&	\hr\cg{(mo-wɛki)}	&	{\hr\hr}lo	&	{\hr\hr}wɛy-ɛ	&	{\hr\hr}mo-tir	&	{\hr\hr}nun\\
E. \ili{Marsela}	&	{\hr\hr}ɛm	&	{\hr\hr}ra	&	\hr\cg{(mo-riay-ei)}	&	{\hr\hr}loi	&	{\hr\hr}uwei	&	{\hr\hr}mo-tir=i	&	{\hr\hr}non\\
\ili{Serili}	&	{\hr\hr}ɛm	&	{\hr\hr}ra-ɛ	&	{\hr\hr}rɛ-way-ɛ	&	{\hr\hr}lɛ	&	{\hr\hr}bia-liɛ	&	{\hr\hr}lya-tir	&	{\hr\hr}non\\
SE. \ili{Babar}	&	{\hr\hr}ɛm	&	{\hr\hr}ra	&	\hr\cg{(mo-yar-e)}	&	{\hr\hr}le	&	{\hr\hr}wiau	&	{\hr\hr}mo-itir	&	{\hr\hr}non\\
\ili{Emplawas}	&	{\hr\hr}im-o	&	{\hr\hr}rɛ-o	&	\hr\cg{(lɛ-wɛt-e)}	&	{\hr\hr}lo	&	{\hr\hr}wɛi-li	&	{\hr\hr}lɛ-tir	&	{\hr\hr}nun-o\\
\ili{Imroing}	&	{\hr\hr}im-e	&	{\hr\hr}reɛ-mo	&	{\hr\hr}mo-wey-e	&	{\hr\hr}lo	&	{\hr\hr}weyeyi	&	{\hr\hr}mo-tiro	&	{\hr\hr}nun-mu\\
Tela	&	{\hr\hr}im-e	&	{\hr\hr}re-nye	&	{\hr\hr}re-wey-e	&	{\hr\hr}no-e	&	{\hr\hr}wei-e	&	{\hr\hr}ne-tir	&	{\hr\hr}nun-mu\\
		\lspbottomrule
	\end{tabular}
		\begin{tablenotes}\tnsize
			\item[†] {W. \ili{Damer} and \tsc{South Babar} reflexes of *diRus
								have C metathesis.
								Initial **wa is probably from *wahiyəR ‘water’.}
		\end{tablenotes}
	\end{threeparttable}
\end{adjustbox}
\end{table}

\begin{table}\centering\tcp{\tsc{Babar} *R/*j > **r}{04-48}
\begin{adjustbox}{width=\textwidth}
	\begin{threeparttable}\st{2pt}
	\begin{tabular}{lll@{}llll}\lsptoprule
local gl.	&	{\hr}tail	&	{\hr}new	&	{\hr}bamboo	&	{\hr}mountain	&	{\hr}nose	&	{\hr}how m.?\\
PMP	&	\hr*iku\tbr{R}	&	\hr*baqə\tbr{R}u	&	\hr*qau\tbr{R}	&	\hr*buki\tbr{j}	&	\hr*i\tbr{j}uŋ	&	\hr*pi\tbr{j}a\\
P.-\ili{Babar}	&	**iu\tbr{r}/**i\tbr{r}u	&	**wa\tbr{r}i\su{†}	&	**o\tbr{r}	&	**wu\tbr{r}\su{‡}	&	**i\tbr{r}n	&	**i\tbr{r}a\\ \midrule
W. \ili{Damer}	&	\hr\cg{(ehɛheya)}	&	\hr\cg{(wɛlwɛka)}	&	{\hr\hr}o\tbr{r}-o	&	{\hr\hr}awo\tbr{r}	&	\hr\cg{(su-no)}	&	\hr\cg{(wi)}\\
Dawera	&	{\hr\hr}i-ol	&	\hr\cg{(wikɛl)}	&	\hr\cg{(welel)}	&	{\hr\hr}wu\tbr{r}	&	\hr\cg{(ud-ol)}	&	{\hr\hr}i\tbr{r}\\
\ili{Dai}	&	{\hr\hr}iʔ-on	&	\hr\cg{(wiwit)}	&	\hr\cg{(fitd-on)}	&	{\hr\hr}we\tbr{r}ʔ-on	&	\hr\cg{(ud-on)}	&	{\hr\hr}i\tbr{r}ə\\
N. \ili{Babar}	&	{\hr\hr}i-ni-ai	&	{\hr\hr}wi{\tl}wiɛ	&	\hr\cg{(spinda)}	&	{\hr\hr}wu\tbr{r}a	&	\hr\cg{(ot-ne)}	&	{\hr\hr}i\tbr{r}e\\
C. \ili{Marsela}	&	{\hr\hr}i\tbr{r}i-nyeo	&	{\hr\hr}war{\tl}wie\tbr{r}i	&	{\hr\hr}o\tbr{r}	&	{\hr\hr}wo\tbr{r}	&	{\hr\hr}i\tbr{r}n-ei	&	{\hr\hr}wie\tbr{r}a\\
W. \ili{Marsela}	&	{\hr\hr}ir{\tl}i\tbr{r}y-ɛn	&	{\hr\hr}war{\tl}wyɛ\tbr{r}	&	{\hr\hr}o\tbr{r}	&	{\hr\hr}wo\tbr{r}	&	{\hr\hr}inn	&	{\hr\hr}wyɛ\tbr{r}a\\
E. \ili{Marsela}	&	{\hr\hr}i-ly-ei	&	{\hr\hr}war{\tl}wia\tbr{r}	&	{\hr\hr}o\tbr{r}	&	{\hr\hr}wo\tbr{r}	&	{\hr\hr}i\tbr{r}n	&	{\hr\hr}wie\tbr{r}a\\
\ili{Serili}	&	{\hr\hr}i-ly-ɛ	&	{\hr\hr}bar{\tl}bia\tbr{r}i	&	{\hr\hr}ɔ\tbr{r}	&	\hr\cg{(bunt-e)}	&	{\hr\hr}i\tbr{r}ɛl	&	\hr\cg{(nomia)}\\
SE. \ili{Babar}	&	{\hr\hr}iy	&	{\hr\hr}wa{\tl}way	&	\hr\cg{(pul)}	&	{\hr\hr}wo\tbr{r}	&	{\hr\hr}i\tbr{r}l	&	\hr\cg{(nomya)}\\
\ili{Emplawas}	&	{\hr\hr}ir{\tl}i\tbr{r}y-o	&	{\hr\hr}war{\tl}wyɛ\tbr{r}y-o	&	{\hr\hr}o\tbr{r}	&	{\hr\hr}wo\tbr{r}	&	\hr(h)itt-ɛ	&	{\hr\hr}wyɛ\tbr{r}a\\
\ili{Imroing}	&	{\hr\hr}ir{\tl}i\tbr{r}i-e	&	\hr\cg{(peruti)}	&	\hr\cg{(lorl-e)}	&	{\hr\hr}wo\tbr{r}u	&	{\hr\hr}i\tbr{r}lɛ-mu	&	{\hr\hr}wo-i\tbr{r}u\\
Tela	&	{\hr\hr}i\tbr{r}y-e	&	{\hr\hr}war{\tl}wiɛ\tbr{r}i	&	\hr\cg{(pɛkk-e)}	&	{\hr\hr}wo\tbr{r}-e	&	{\hr\hr}i\tbr{r}n-e	&	{\hr\hr}wɔ-i\tbr{r}i\\
		\lspbottomrule
	\end{tabular}
		\begin{tablenotes}\tnsize
			\item[†]	{Reflexes are reduplicated forms
								with metathesis of *i across a morpheme boundary.}
			\item[‡]	{The reflexes in Dawera and
								North \ili{Babar} are given for both ‘mountain’ and ‘forest’.
								The reflexes in Tela
								and \ili{Imroing} are given for ‘forest’.
								\ili{Emplawas} \it{wor} is ‘forest’.}
		\end{tablenotes}
	\end{threeparttable}
\end{adjustbox}
\end{table}

\largerpage
It is unlikely that \tsc{Babar} *z > **h was a direct change.
\ili{As} discussed at the beginning of \srf{04-06},
it probably went through intermediate **s.
\tsc{\ili{Teun}-\ili{Nila}-Serua} also has *z > \it{s}.
However, we cannot place \tsc{\ili{Teun}-\ili{Nila}-Serua}
and \tsc{Babar} together on the basis of shared *z > **s because
\tsc{Babar} has *s > **ʈ which does not merge with *z.
Thus, putative *z > **s must have occurred in \tsc{Babar}
after *s > **ʈ and thus after its separation 
from other languages of \tsc{Southwest Maluku}.

\begin{table}\centering\tcp{\tsc{Babar} *w/*b > **w}{BabWB}
\begin{adjustbox}{width=\textwidth}
	%\begin{threeparttable}
	\st{2pt}\begin{tabular}{lllllll}\lsptoprule
local gl.	&	(fresh) water	&	ySib	&	root	&	moon	&	fruit	&	dust, ash	\\	
PMP	&	\hr*\tbr{w}ahiyəR	&	\hr*huaji	&	\hr*\tbr{w}akaR	&	\hr*\tbr{b}atu	&	\hr*\tbr{b}uaq	&	\hr*qa\tbr{b}u	\\	
P.-\ili{Babar}	&	**\tbr{w}ey	&	**\tbr{w}ayi	&	**\tbr{w}iay	&	**\tbr{w}atu	&	**\tbr{w}u-	&	**(li)-a\tbr{w}u	\\	\midrule
W. \ili{Damer}	&	{\hr\hr}\tbr{w}i-o	&	{\hr\hr}\tbr{w}es-eni	&	{\hr\hr}\tbr{w}oto	&	{\hr\hr}\tbr{w}oth-o	&	{\hr\hr}uh-a	&	{\hr\hr}a\tbr{w}or	\\	
Dawera	&	{\hr\hr}\tbr{w}ee	&	{\hr\hr}\tbr{w}el-ol	&	{\hr\hr}\tbr{w}oy-el-ol	&	{\hr\hr}\tbr{w}atk	&	{\hr\hr}uk	&	{\hr\hr}le\tbr{w}k-ol	\\	
\ili{Dai}	&	{\hr\hr}\tbr{w}ɛy-on	&	{\hr\hr}\tbr{w}ai-nyon	&	{\hr\hr}\tbr{o}i-on	&	{\hr\hr}\tbr{w}atʔ-on	&	{\hr\hr}\tbr{w}uʔ-ɛn-on	&	{\hr\hr}nya\tbr{w}ʔ-on	\\	
N. \ili{Babar}	&	{\hr\hr}\tbr{w}iey-a	&	{\hr\hr}\tbr{w}ia-nyai	&	{\hr\hr}\tbr{w}ia-na	&	{\hr\hr}\tbr{w}aʔay-ai	&	{\hr\hr}\tbr{w}una	&	{\hr\hr}nia\tbr{w}i	\\	
C. \ili{Marsela}	&	{\hr\hr}\tbr{w}ee	&	{\hr\hr}\tbr{w}ei-ni	&	{\hr\hr}\tbr{w}ɛ-nye	&	{\hr\hr}\tbr{w}ay	&	{\hr\hr}\tbb{b}ui-ni, \tbr{w}ui-ni	&	{\hr\hr}ɛ\tbr{w}y	\\	
W. \ili{Marsela}	&	{\hr\hr}\tbr{w}e	&	{\hr\hr}\tbr{w}ɛy-ɛ	&	{\hr\hr}\tbr{w}ɛ-ni	&	{\hr\hr}\tbr{w}aky	&	{\hr\hr}\tbr{w}uy	&	{\hr\hr}u\tbr{w}y	\\	
E. \ili{Marsela}	&	{\hr\hr}\tbr{w}e	&	{\hr\hr}\tbr{uw}ei	&	{\hr\hr}i\tbr{w}ier-ne	&	{\hr\hr}\tbr{w}ak	&	{\hr\hr}\tbr{w}u-liei	&	{\hr\hr}u\tbr{w}	\\	
\ili{Serili}	&	{\hr\hr}\tbb{b}e	&	{\hr\hr}\tbb{b}ia-liɛ	&	{\hr\hr}\tbb{b}ia-le	&	{\hr\hr}\tbb{b}aʔay-e	&	{\hr\hr}\tbb{b}u-lye	&	{\hr\hr}u\tbb{b}i-e, u	\\	
SE. \ili{Babar}	&	{\hr\hr}\tbr{w}ey	&	{\hr\hr}\tbr{w}iau	&	{\hr\hr}\tbr{w}ia	&	{\hr\hr}\tbr{w}axay	&	{\hr\hr}\tbr{w}u	&	{\hr\hr}u\tbr{w}y, u\tbb{b}i	\\	
\ili{Emplawas}	&	{\hr\hr}\tbr{w}e	&	{\hr\hr}\tbr{w}ɛi-li	&	{\hr\hr}\tbr{w}ɛ-lyo	&	{\hr\hr}\tbr{w}aky-o	&	{\hr\hr}\tbr{w}ui-li	&	{\hr\hr}u\tbr{w}y-o	\\	
\ili{Imroing}	&	{\hr\hr}\tbr{w}ey-e	&	{\hr\hr}\tbr{w}eyeyi	&	{\hr\hr}\tbr{w}ɛ-lye	&	{\hr\hr}\tbr{w}aky-e	&	{\hr\hr}\tbr{w}ui-lyɛ	&	{\hr\hr}u\tbr{w}i-e	\\	
Tela	&	{\hr\hr}\tbr{w}ey-e	&	{\hr\hr}\tbr{w}ei-e	&	{\hr\hr}\tbr{w}ɛ-niə	&	{\hr\hr}\tbr{w}ak-e	&	{\hr\hr}\tbr{w}u-nyɛ	&	{\hr\hr}o\tbb{b}i-e	\\	
		\lspbottomrule
	\end{tabular}
		%\begin{tablenotes}\tnsize
			%\item[†]
			%%\item[‡]
			%%\item[Δ]
			%%\item[‖]
			%%\item[¶]
		%\end{tablenotes}
	%\end{threeparttable}
\end{adjustbox}
\end{table}

Like nearly all other \tsc{Southwest Maluku} languages,
*w/*b usually merge in \tsc{Babar}.
With the exception of \ili{Serili}, the outcome of this
merger in the data in \citet{Taber1993}
is a labial and/or labio-dental fricative
([β], [v]) and/or approximant ([w], [ʋ]).
Examples are given in \trf{BabWB},
where all these realisations are transcribed ‹w›.

\ili{Serili} has *w/*b > \it{b} and it this could be due to *w > \it{b}
with retention of *b unchanged.
However, other \tsc{Babar} languages also have sporadic instances of *b = \it{b}.
In \ili{Emplawas} -- for which the most phonetic data is
available -- /w/ is usually a bilabial
approximant [β̞] or labio-velar approximant [w].
Adjacent to the high-back vowel /u/ it is often
a bilabial fricative [β] and, occasionally, a plosive [b].
Given all this data, we prefer to analyse \ili{Serili} *w/*b > \it{b}
as due to fortition of intermediate **w (perhaps [β])
rather than due to retention of *b unchanged;
thus we propose \ili{Serili} *w/*b > **w > \it{b}.
